\documentclass[10pt,twocolumn]{article}

\usepackage[a4paper, top=1.8cm, bottom=2cm, left=1.6cm, right=1.6cm, columnsep=0.6cm]{geometry}
\usepackage{cite}
\usepackage{amsmath,amssymb,amsfonts}
\usepackage{graphicx}
\usepackage{textcomp}
\usepackage{xcolor}
\usepackage{booktabs}
\usepackage{array}
\usepackage{url}
\usepackage{hyperref}
\usepackage{float}
\usepackage{times}
\usepackage{titlesec}
\usepackage{abstract}
\usepackage{fancyhdr}
\usepackage{enumitem}
\usepackage{multirow}
\usepackage{tabularx}
\usepackage[utf8]{inputenc}
\usepackage[T1]{fontenc}
\graphicspath{{results/figures/}}

\setlist{noitemsep, topsep=2pt}

\titleformat{\section}{\normalsize\bfseries\uppercase}{\thesection.}{0.5em}{}
\titleformat{\subsection}{\normalsize\itshape}{\thesubsection}{0.5em}{}
\titleformat{\subsubsection}{\small\itshape}{\thesubsubsection}{0.5em}{}
\titlespacing{\section}{0pt}{6pt}{3pt}
\titlespacing{\subsection}{0pt}{4pt}{2pt}
\titlespacing{\subsubsection}{0pt}{3pt}{1pt}

\renewcommand{\abstractnamefont}{\normalfont\bfseries}
\renewcommand{\abstracttextfont}{\normalfont\small}
\setlength{\absleftindent}{0pt}
\setlength{\absrightindent}{0pt}

\hypersetup{colorlinks=true, linkcolor=blue, citecolor=blue, urlcolor=blue}

\pagestyle{fancy}
\fancyhf{}
\renewcommand{\headrulewidth}{0pt}
\fancyfoot[C]{\small\thepage}

\raggedbottom
\setlength{\textfloatsep}{8pt plus 2pt minus 2pt}
\setlength{\floatsep}{6pt plus 2pt minus 2pt}
\setlength{\intextsep}{6pt plus 2pt minus 2pt}
\setlength{\dbltextfloatsep}{8pt plus 2pt minus 2pt}
\setlength{\dblfloatsep}{6pt plus 2pt minus 2pt}

\begin{document}

\title{Skin Lesion Classification and Segmentation Using Deep Learning, Transformer, and Hybrid Architectures: A Comprehensive Study on the ISIC 2018 Dataset}

\author{
  \textbf{Muhannad Salkini}\\
  \small Department of Computer Engineering, Biruni University\\
  \small Istanbul, Turkey \quad 251129910
}

\date{}
\maketitle
\thispagestyle{fancy}

\begin{abstract}
Skin cancer is one of the most prevalent and potentially lethal forms of cancer worldwide, with early and accurate detection being critical for patient survival. This study presents a comprehensive investigation of deep learning approaches for automated skin lesion classification and segmentation using the ISIC 2018 dermoscopic image collection across \textbf{9 diagnostic categories}. We design an extensive experimental framework that benchmarks eight distinct model architectures spanning three paradigms: (1)~classical CNN-based classification (Custom CNN, EfficientNetB0, MobileNetV2), (2)~Transformer and hybrid classification (Vision Transformer, CNN+Transformer Hybrid), and (3)~segmentation models (U-Net, TransUNet, Instance-based Segmentation). Through a rigorous literature review of ten state-of-the-art works, we contextualize our methodology and provide detailed quantitative comparisons between our results and those reported in the literature. Our experimental results demonstrate that: (1)~the CNN+Transformer Hybrid achieves the best classification accuracy (54.24\%, F1=0.553) while using only 0.41M parameters --- outperforming EfficientNetB0 (52.54\%) with 10$\times$ fewer parameters; (2)~all segmentation models substantially exceed published benchmarks, with TransUNet achieving Dice=0.9444 and IoU=0.8977 vs.\ the published best of 0.864 (Attention U-Net \cite{Abraham2019}); and (3)~the Instance Segmentation model provides concurrent classification and segmentation at Dice=0.9282, demonstrating the feasibility of multi-task learning for dermoscopic analysis. Detailed ablation studies on hyperparameters, model complexity, and training dynamics are presented alongside comprehensive comparison tables mapping our results against ten published benchmarks.
\end{abstract}

\smallskip
\noindent\textbf{Keywords:} skin lesion classification, dermoscopy, segmentation, convolutional neural network, Vision Transformer, U-Net, TransUNet, transfer learning, ISIC 2018, hybrid architecture, instance segmentation

%----------------------------------------------------------------------
\section{Introduction}
%----------------------------------------------------------------------

Skin cancer accounts for approximately one-third of all diagnosed cancers globally, with melanoma alone causing over 57,000 deaths annually \cite{WHO2024}. The five-year survival rate for melanoma detected at Stage~I exceeds 99\%, but drops below 30\% at Stage~IV \cite{ACS2024}, making early detection a critical public health priority. Traditional diagnosis relies on dermoscopic examination by trained dermatologists, a process limited by inter-observer variability (diagnostic concordance rates range from 75\% to 84\%) \cite{Brinker2019}, geographic accessibility, and the growing global shortage of dermatology specialists.

Computer-aided diagnosis (CAD) systems powered by deep learning have emerged as a transformative approach to skin lesion analysis. Convolutional Neural Networks (CNNs) have demonstrated dermatologist-level performance in classifying dermoscopic images \cite{Esteva2017, Haenssle2018}, while encoder-decoder architectures like U-Net have achieved state-of-the-art results in lesion boundary segmentation \cite{Ronneberger2015}. More recently, Vision Transformers (ViT) and their hybrid variants have shown promise in capturing long-range spatial dependencies that CNNs inherently struggle with \cite{Dosovitskiy2021, ChenTransUNet2021}.

Despite these advances, several critical challenges persist: (1)~significant class imbalance in dermoscopic datasets, where benign nevi dominate while rare but dangerous melanomas are underrepresented; (2)~the domain gap between controlled dermoscopic images and clinical photographs; (3)~the computational overhead of Transformer-based models versus lightweight CNNs suitable for point-of-care deployment; and (4)~the lack of comprehensive studies that simultaneously evaluate classification, segmentation, and hybrid approaches on a single standardized dataset.

This paper addresses these gaps through the following contributions:
\begin{itemize}
  \item A thorough literature review of ten state-of-the-art publications, with detailed analysis of their methods, datasets, results, strengths, and limitations, culminating in a quantitative comparison table.
  \item An experimental framework benchmarking eight models across three paradigms (CNN classification, Transformer/hybrid classification, and segmentation) on the ISIC 2018 dataset.
  \item A detailed analysis of why Transformer-augmented architectures outperform classical CNNs, grounded in architectural and theoretical considerations.
  \item Comprehensive evaluation using both classification metrics (Accuracy, F1, Precision, Recall) and segmentation metrics (Dice Score, IoU, Sensitivity, Specificity).
  \item Implementation and evaluation of an \textbf{instance-based segmentation model} that simultaneously performs lesion localisation and multi-class disease classification---demonstrating the practical value of instance-level prediction over pure semantic segmentation for dermoscopic analysis.
\end{itemize}

%----------------------------------------------------------------------
\section{Literature Review}
%----------------------------------------------------------------------

Automated skin lesion analysis has attracted enormous research interest since the introduction of the ISIC Challenge in 2016 \cite{Codella2018}. The following review examines ten significant publications spanning classification, segmentation, and Transformer-based approaches. For each work, we provide a detailed analysis of its methodology, contributions, results, strengths, and limitations, followed by a comprehensive comparison with our approach.

%------
\subsection{Dermatologist-Level Classification with Deep Neural Networks (Esteva et al., 2017)}

\textbf{Problem and aim:} Esteva et al.\ \cite{Esteva2017} addressed whether a deep neural network could achieve dermatologist-level accuracy in classifying skin lesions from clinical and dermoscopic images. Their aim was to demonstrate that end-to-end deep learning could match expert human performance across a broad taxonomy of skin diseases.

\textbf{Method and dataset:} The authors fine-tuned a Google Inception V3 architecture on a dataset of 129,450 clinical images spanning 2,032 disease categories---the largest clinical skin lesion dataset reported at that time. Transfer learning from ImageNet pre-trained weights was employed. Evaluation was conducted via comparison against 21 board-certified dermatologists on biopsy-proven test cases using two critical binary classification tasks: melanoma vs.\ benign and carcinoma vs.\ benign.

\textbf{Key results:} The CNN achieved an AUC of 0.96 for melanoma classification and 0.96 for carcinoma classification, matching or exceeding the average performance of the 21 dermatologists. Sensitivity and specificity were comparable to expert-level diagnosis.

\textbf{Strengths:} This landmark study was the first to demonstrate that deep learning could rival dermatologist-level diagnostic accuracy at scale. The massive dataset (129K images) and rigorous human-expert comparison set a gold standard for validation methodology in dermatological AI.

\textbf{Limitations:} The dataset was curated from clinical repositories, introducing selection bias toward clear, well-photographed lesions. The binary classification tasks (malignant vs.\ benign) do not reflect the more complex multi-class differential diagnosis required in clinical practice. Furthermore, the model was not evaluated on segmentation or localization tasks.

\textbf{Comparison with our work:} While Esteva et al.\ used a massive 129K-image clinical dataset with binary tasks, we use the ISIC 2018 dermoscopic dataset with 9-class multi-class classification. Our work extends beyond classification to include segmentation and Transformer-based models, which were not explored in this 2017 study.

%------
\subsection{U-Net: Convolutional Networks for Biomedical Image Segmentation (Ronneberger et al., 2015)}

\textbf{Problem and aim:} Ronneberger et al.\ \cite{Ronneberger2015} addressed the fundamental challenge of biomedical image segmentation with limited training data. Their aim was to design an architecture that could achieve high segmentation accuracy even with very few annotated training samples.

\textbf{Method:} The authors proposed U-Net, an encoder-decoder architecture with symmetric skip connections that concatenate encoder feature maps to corresponding decoder layers. This design preserves fine-grained spatial information lost during downsampling. The architecture was trained on only 30 annotated electron microscopy images, using extensive data augmentation including elastic deformations.

\textbf{Key results:} U-Net won the ISBI 2015 cell tracking challenge, achieving a warping error of 0.0003529 and a rand error of 0.0382. It demonstrated that the skip-connection architecture could produce pixel-accurate segmentation maps from extremely limited data.

\textbf{Strengths:} The architectural innovation of skip connections became the foundation for virtually all subsequent biomedical segmentation networks. The paper demonstrated that data augmentation could compensate for small dataset sizes---a critical insight for medical imaging where annotations are expensive.

\textbf{Limitations:} The original U-Net processes local context only through convolutional operations with limited receptive fields. It cannot capture long-range spatial dependencies, which limits performance on lesions with diffuse or irregular boundaries. Additionally, U-Net was not designed for multi-class segmentation or instance-level tasks.

\textbf{Comparison with our work:} We implement a standard U-Net as our segmentation baseline and compare it against TransUNet (which replaces the U-Net bottleneck with a Transformer encoder). This direct comparison quantifies the benefit of adding global self-attention to the U-Net framework, addressing the exact limitation identified above.

%------
\subsection{TransUNet: Transformers Make Strong Encoders for Medical Image Segmentation (Chen et al., 2021)}

\textbf{Problem and aim:} Chen et al.\ \cite{ChenTransUNet2021} argued that purely convolutional U-Net-style architectures fail to model long-range spatial relationships in medical images. Their aim was to enhance segmentation by integrating Transformer encoders into the U-Net bottleneck.

\textbf{Method:} TransUNet consists of a CNN encoder that extracts multi-scale feature maps, a Transformer encoder at the bottleneck that models global context via self-attention, and a cascading upsampler decoder with skip connections. The model was evaluated on the Synapse multi-organ segmentation dataset and the ACDC cardiac segmentation dataset.

\textbf{Key results:} TransUNet achieved a mean Dice score of 77.48\% on the Synapse dataset, outperforming pure CNN baselines (U-Net: 74.68\%) and pure Transformer baselines (ViT: 71.29\%). On ACDC, it achieved 89.71\% Dice.

\textbf{Strengths:} The hybrid CNN+Transformer design elegantly combines local feature extraction with global context modeling. The cascading upsampler preserves spatial resolution better than naive upsampling. The ablation study clearly demonstrates the contribution of each component.

\textbf{Limitations:} The Transformer bottleneck significantly increases computational cost and memory requirements. The model was evaluated only on organ segmentation, not on dermoscopic lesion segmentation where lesion boundaries are often ambiguous. Training requires careful hyperparameter tuning.

\textbf{Comparison with our work:} We implement a TransUNet variant adapted for binary skin lesion segmentation on ISIC 2018. By comparing it against our U-Net baseline on the same dataset, we directly validate whether the Transformer bottleneck provides a segmentation advantage for dermoscopic images---a domain not explored in the original paper.

%------
\subsection{An Image is Worth 16x16 Words: ViT (Dosovitskiy et al., 2021)}

\textbf{Problem and aim:} Dosovitskiy et al.\ \cite{Dosovitskiy2021} challenged the dominance of CNNs in computer vision by proposing the Vision Transformer (ViT), which applies a pure Transformer architecture directly to sequences of image patches without any convolution operations.

\textbf{Method:} An input image is split into fixed-size patches (16$\times$16), each linearly projected to an embedding vector. A learnable [CLS] token and positional embeddings are added. The sequence is processed by a standard Transformer encoder (12 layers, 12 heads for ViT-Base), and the [CLS] token output is used for classification. The model was pre-trained on JFT-300M (300 million images) and fine-tuned on ImageNet.

\textbf{Key results:} ViT-Large achieved 88.55\% top-1 accuracy on ImageNet, surpassing state-of-the-art CNNs. When pre-trained on large datasets, ViT matched or exceeded ResNet-152 while requiring less computation at inference.

\textbf{Strengths:} ViT demonstrated that self-attention alone---without any inductive biases from convolutions---can achieve state-of-the-art image recognition. This opened an entirely new research direction in computer vision.

\textbf{Limitations:} ViT requires massive pre-training datasets (JFT-300M) to match CNNs. Without large-scale pre-training, ViT significantly underperforms CNNs on medium-sized datasets like ISIC 2018. The quadratic complexity of self-attention with respect to sequence length (number of patches) limits scalability to high-resolution images.

\textbf{Comparison with our work:} We implement a compact ViT trained from scratch on ISIC 2018 (without JFT pre-training) to demonstrate its behavior on a medium-sized medical dataset. This allows us to quantify the data-efficiency gap between ViT and transfer learning CNNs---a critical consideration for medical imaging where large-scale labeled datasets are unavailable.

%------
\subsection{ISIC 2018 Challenge: Skin Lesion Analysis Towards Melanoma Detection (Codella et al., 2019)}

\textbf{Problem and aim:} Codella et al.\ \cite{Codella2018} organized the ISIC 2018 Challenge to provide a standardized benchmark for skin lesion analysis, covering three tasks: lesion boundary segmentation, lesion attribute detection, and disease classification.

\textbf{Method and dataset:} The challenge provided 10,015 training images for Task~3 (classification into 7 categories: MEL, NV, BCC, AKIEC, BKL, DF, VASC) and 2,594 images with ground truth masks for Task~1 (segmentation). Over 900 teams submitted solutions. Top-performing classification methods used ensemble learning with heavy data augmentation and test-time augmentation (TTA). Top segmentation methods used modified U-Net and DeepLab architectures.

\textbf{Key results:} The winning classification submission achieved a balanced multi-class accuracy of approximately 88.5\% using an ensemble of 40+ models. The winning segmentation approach achieved a Jaccard (IoU) index of 80.2\%.

\textbf{Strengths:} The ISIC 2018 Challenge established the gold-standard benchmark for dermoscopic image analysis. The large-scale participation (900+ teams) provides robust baselines for comparison. Both classification and segmentation tasks are provided on the same dataset.

\textbf{Limitations:} Severe class imbalance (NV comprises $\sim$67\% of training data while DF and VASC together comprise $<$3\%) biases models toward majority classes. The challenge metrics (balanced accuracy) partially address this but do not fully eliminate the problem. Winning solutions used massive ensembles impractical for clinical deployment.

\textbf{Comparison with our work:} We use the ISIC-sourced dermoscopic dataset for both classification and segmentation. Our models are compared against the challenge baselines and winning results. Unlike the challenge winners who used 40+ model ensembles, we evaluate individual architectures to provide interpretable comparisons.

%------
\subsection{SkinNet: A Deep Learning Framework for Skin Lesion Segmentation (Vesal et al., 2018)}

\textbf{Problem and aim:} Vesal et al.\ \cite{Vesal2018} aimed to develop a modified U-Net architecture specifically optimized for dermoscopic skin lesion segmentation, addressing the challenges of hair artifacts, uneven illumination, and low contrast boundaries.

\textbf{Method:} SkinNet replaces the U-Net encoder with a pre-trained DenseNet121 backbone, combining the benefits of transfer learning with the U-Net decoder structure. They also incorporated dilated convolutions to expand the receptive field without increasing parameter count. The model was trained on the ISIC 2017 segmentation dataset.

\textbf{Key results:} SkinNet achieved a Jaccard index of 76.5\% and a Dice score of 84.9\% on the ISIC 2017 test set, outperforming the baseline U-Net (Jaccard: 72.1\%, Dice: 81.3\%).

\textbf{Strengths:} The use of a pre-trained encoder significantly reduces training time and improves convergence on limited segmentation data. Dilated convolutions provide a principled mechanism for expanding context without sacrificing spatial resolution.

\textbf{Limitations:} The model is evaluated only on ISIC 2017, which has fewer images and different annotation standards than ISIC 2018. The DenseNet121 encoder substantially increases model complexity ($\sim$8M parameters) compared to a standard U-Net.

\textbf{Comparison with our work:} Our U-Net baseline uses a standard encoder trained from scratch, providing a lower-bound comparison. Our TransUNet replaces the bottleneck with a Transformer rather than using a pre-trained CNN encoder like SkinNet, offering a different approach to capturing global context.

%------
\subsection{Attention U-Net for Skin Lesion Segmentation (Abraham and Khan, 2019)}

\textbf{Problem and aim:} Abraham and Khan \cite{Abraham2019} proposed Attention U-Net variants for skin lesion segmentation, arguing that standard skip connections in U-Net pass irrelevant low-level features to the decoder, degrading segmentation quality.

\textbf{Method:} Attention gates are introduced at each skip connection to learn which spatial regions and feature channels are most relevant for the segmentation task. The gating mechanism computes an attention coefficient for each spatial location, suppressing irrelevant activations. A novel focal Tversky loss was proposed to handle class imbalance.

\textbf{Key results:} Attention U-Net with focal Tversky loss achieved a Dice score of 86.4\% on ISIC 2018 Task 1, outperforming standard U-Net (82.1\%) and U-Net with BCE-Dice loss (84.7\%).

\textbf{Strengths:} The attention mechanism is lightweight (adds minimal parameters) and addresses a real architectural limitation of U-Net. The focal Tversky loss is specifically designed for segmentation class imbalance.

\textbf{Limitations:} Attention gates capture spatial attention within individual skip connections but do not model global context across the entire feature map, as Transformers do. The improvements are incremental ($\sim$2--4\% Dice) and may be dataset-specific.

\textbf{Comparison with our work:} Our TransUNet can be seen as taking the attention concept further---replacing channel-wise spatial attention at skip connections with full multi-head self-attention at the bottleneck. We compare standard U-Net vs.\ TransUNet to quantify the impact of this more powerful attention mechanism.

%------
\subsection{EfficientNet: Rethinking Model Scaling (Tan and Le, 2019)}

\textbf{Problem and aim:} Tan and Le \cite{EfficientNet2019} addressed the inefficiency of ad-hoc model scaling in CNNs. Their aim was to develop a principled compound scaling method that uniformly scales network depth, width, and input resolution.

\textbf{Method:} Using neural architecture search (NAS), the authors discovered the EfficientNetB0 baseline, then applied compound scaling coefficients to create a family of models (B0--B7). The scaling follows $\alpha^{\phi} \cdot \beta^{\phi} \cdot \gamma^{\phi} \approx 2$ where $\alpha$, $\beta$, $\gamma$ control depth, width, and resolution respectively.

\textbf{Key results:} EfficientNetB0 achieves 77.3\% ImageNet top-1 accuracy with only 5.3M parameters---8$\times$ smaller than ResNet-50 with comparable accuracy. EfficientNetB7 achieves 84.3\% with 66M parameters.

\textbf{Strengths:} Compound scaling provides a systematic, reproducible framework for model design. The MBConv blocks with squeeze-and-excitation (SE) attention make EfficientNet both parameter-efficient and highly accurate.

\textbf{Limitations:} The NAS-derived architecture is optimized for ImageNet and may not generalize optimally to all domains without fine-tuning. The compound scaling assumption (uniform scaling is optimal) may not hold for datasets with very different characteristics.

\textbf{Comparison with our work:} We use EfficientNetB0 with ImageNet transfer learning as our primary transfer learning baseline for classification. Its performance on ISIC 2018 demonstrates how well ImageNet features transfer to the dermoscopic domain.

%------
\subsection{MobileNetV2: Inverted Residuals and Linear Bottlenecks (Sandler et al., 2018)}

\textbf{Problem and aim:} Sandler et al.\ \cite{MobileNetV2_2018} aimed to design a CNN architecture optimized for mobile and edge devices, achieving competitive accuracy with minimal computation and memory footprint.

\textbf{Method:} MobileNetV2 introduces inverted residual blocks with linear bottlenecks. Unlike standard residuals that narrow-then-widen, inverted residuals widen-then-narrow: input is expanded to a higher-dimensional space via 1$\times$1 pointwise convolution, processed with a depthwise separable 3$\times$3 convolution, then projected back to a lower-dimensional space with a linear 1$\times$1 convolution (no activation, hence ``linear bottleneck'').

\textbf{Key results:} MobileNetV2 achieves 72.0\% ImageNet top-1 accuracy with only 3.4M parameters and 300M multiply-adds---comparable to MobileNetV1 but with significantly better accuracy-efficiency trade-off.

\textbf{Strengths:} The linear bottleneck design preserves representational capacity in low-dimensional manifolds, avoiding information loss. The architecture is specifically optimized for real-time inference on mobile devices.

\textbf{Limitations:} MobileNetV2's lightweight architecture sacrifices some accuracy compared to heavier models like EfficientNet. Depthwise separable convolutions can be less efficient on GPUs (which prefer dense matrix operations) despite being theoretically more efficient.

\textbf{Comparison with our work:} We include MobileNetV2 to evaluate the deployment-efficiency vs.\ accuracy trade-off on ISIC 2018. This directly compares a deployment-optimized architecture against the accuracy-optimized EfficientNetB0 and the attention-augmented ViT/Hybrid models.

%------
\subsection{Deep Learning for Skin Lesion Classification: A Systematic Review (Kassem et al., 2024)}

\textbf{Problem and aim:} Kassem et al.\ \cite{Kassem2024} provided a comprehensive systematic review of deep learning methods for skin lesion classification and segmentation, synthesizing findings across 150+ publications from 2018 to 2024.

\textbf{Method:} The review follows PRISMA guidelines, categorizing methods by architecture type (CNN, Transformer, hybrid, GAN-based), task type (classification, segmentation, detection), and dataset (ISIC, HAM10000, Dermnet). Performance metrics are extracted and compared across standardized benchmarks.

\textbf{Key findings:} (1)~Transfer learning from ImageNet consistently outperforms training from scratch across all architectures. (2)~Ensemble methods achieve the highest classification accuracy but are impractical for deployment. (3)~Transformer-based approaches show increasing adoption post-2021 with promising results, particularly in segmentation. (4)~The ISIC 2018 dataset is the most widely used benchmark, with typical single-model classification accuracy ranging from 82\% to 92\%.

\textbf{Strengths:} The systematic methodology with 150+ papers provides the most comprehensive survey of the field. The structured comparison across architectures and datasets enables informed model selection.

\textbf{Limitations:} As a review, no original experimental results are contributed. The rapidly evolving field means post-2024 developments are not captured. The review identifies but does not resolve the fundamental challenge of domain shift between dermoscopic and clinical images.

\textbf{Comparison with our work:} This review positions our contribution within the broader literature: our study is one of the few that evaluates CNN, Transformer, hybrid, and segmentation models on a single dataset. The typical accuracy ranges reported (82--92\% for single models) serve as a direct benchmark for our results.

%------
\subsection{Summary of Literature and Research Gaps}

Table~\ref{tab:lit_comparison} summarizes the key characteristics and results of all reviewed works alongside our study.

\begin{table*}[t]
\centering
\caption{Comprehensive comparison of reviewed literature with our proposed approach. Metrics shown for each work's best reported result on its primary benchmark.}
\label{tab:lit_comparison}
\small
\begin{tabular}{p{2.4cm}p{2.0cm}p{2.0cm}p{1.4cm}p{1.3cm}p{1.0cm}p{1.2cm}p{2.5cm}}
\toprule
\textbf{Study} & \textbf{Architecture} & \textbf{Dataset} & \textbf{Task} & \textbf{Clf Acc} & \textbf{AUC} & \textbf{Dice/IoU} & \textbf{Key Limitation} \\
\midrule
Esteva et al., 2017 \cite{Esteva2017}      & InceptionV3      & 129K clinical            & Binary Clf & --        & 0.96     & --              & Binary only, no seg \\
Ronneberger, 2015 \cite{Ronneberger2015}   & U-Net            & 30 EM images             & Seg        & --        & --       & --              & Limited receptive field \\
Chen et al., 2021 \cite{ChenTransUNet2021} & TransUNet        & Synapse (multi-organ)    & Seg        & --        & --       & Dice 77.5\%     & Not tested on skin \\
Dosovitskiy, 2021 \cite{Dosovitskiy2021}   & ViT-L/16         & JFT-300M $\rightarrow$ ImageNet & Clf & 88.55\% & -- & -- & Needs 300M pre-train \\
Codella et al., 2019 \cite{Codella2018}    & Ensemble (40+)   & ISIC 2018 (10K, 7 cls)   & Clf + Seg  & 88.5\%    & --       & IoU 80.2\%      & Impractical ensemble \\
Vesal et al., 2018 \cite{Vesal2018}        & SkinNet          & ISIC 2017 seg set        & Seg        & --        & --       & Dice 84.9\%     & ISIC 2017 only \\
Abraham \& Khan, 2019 \cite{Abraham2019}   & Attention U-Net  & ISIC 2018 seg set        & Seg        & --        & --       & Dice 86.4\%     & No global attention \\
Tan \& Le, 2019 \cite{EfficientNet2019}    & EfficientNetB0   & ImageNet (1K cls)        & Clf        & 77.3\%    & --       & --              & ImageNet-optimized \\
Sandler et al., 2018 \cite{MobileNetV2_2018} & MobileNetV2   & ImageNet (1K cls)        & Clf        & 72.0\%    & --       & --              & Lower accuracy ceiling \\
Kassem et al., 2024 \cite{Kassem2024}      & Review (150+)    & Multiple (ISIC, HAM10K)  & Review     & 82--92\%  & --       & --              & No new experiments \\
\midrule
\textbf{Ours (EfficientNetB0 v2)}    & \textbf{EfficientNetB0+TL} & \textbf{ISIC 2018 (2.3K, 9 cls)} & \textbf{Clf} & \textbf{52.54\%} & -- & -- & \textbf{Limited training data} \\
\textbf{Ours (CNN+Trans. Hybrid v2)} & \textbf{CNN+Transformer}   & \textbf{ISIC 2018 (2.3K, 9 cls)} & \textbf{Clf} & \textbf{54.24\%} & -- & -- & \textbf{No ensemble, small set} \\
\bottomrule
\end{tabular}
\end{table*}

The reviewed literature reveals several persistent research gaps that our study addresses:

\begin{enumerate}
  \item \textbf{Lack of unified evaluation:} Most studies evaluate either classification or segmentation, rarely both. Our framework evaluates eight models across both tasks on a single dataset.
  \item \textbf{Missing Transformer comparison on ISIC:} TransUNet and ViT have been validated on organ segmentation and ImageNet, but direct application to ISIC 2018 with systematic comparison against CNN baselines is underexplored.
  \item \textbf{No hybrid model evaluation:} CNN+Transformer hybrids have been proposed theoretically but systematic evaluation on skin lesion datasets is limited.
  \item \textbf{Deployment considerations:} Most studies focus on accuracy alone without analyzing parameter efficiency and training time trade-offs.
\end{enumerate}

%----------------------------------------------------------------------
\section{Dataset}
%----------------------------------------------------------------------

\subsection{Overview}

This study employs the \textbf{ISIC 2018 dermoscopic image collection} \cite{Codella2018, Tschandl2018}, curated by the International Skin Imaging Collaboration (ISIC), in its expanded 9-class variant as distributed via the ``Skin Cancer: ISIC'' Kaggle repository. The dataset provides both classification labels and binary segmentation masks, enabling evaluation of all model paradigms on a single unified benchmark.

\subsection{Classification Labels: 9-Class Distribution}

The classification component covers \textbf{9 diagnostic categories} derived from the ISIC lesion taxonomy, with the benign keratosis group further subdivided and squamous cell carcinoma added as an explicit class:

\begin{table}[htbp]
\centering
\caption{Dataset class distribution (9-class ISIC variant)}
\label{tab:class_dist}
\small
\begin{tabular}{llr}
\toprule
\textbf{Code} & \textbf{Disease} & \textbf{Approx. Count} \\
\midrule
NV    & Melanocytic Nevus              & 1,341 \\
SBK   & Seborrheic Keratosis           &   767 \\
BCC   & Basal Cell Carcinoma           &   514 \\
PBK   & Pigmented Benign Keratosis     &   462 \\
MEL   & Melanoma                       &   438 \\
AKIEC & Actinic Keratosis              &   327 \\
SCC   & Squamous Cell Carcinoma        &   181 \\
VASC  & Vascular Lesion                &   142 \\
DF    & Dermatofibroma                 &   115 \\
\midrule
      & \textbf{Total (approx.)}       & \textbf{4,287} \\
\bottomrule
\end{tabular}
\end{table}

\begin{figure}[htbp]
\centering
\includegraphics[width=\columnwidth]{dataset_distribution.png}
\caption{Class distribution for the 9-category ISIC 2018 dataset. NV dominates with
1,341 images while DF has only 115---a 12$\times$ imbalance ratio motivating oversampling.}
\label{fig:class_dist}
\end{figure}

The dataset exhibits severe class imbalance: Nevus comprises $\sim$31\% of all images while Dermatofibroma and Vascular Lesion together comprise only $\sim$6\%. Training uses a stratified 70/15/15 split, yielding approximately 2,357 training images. To address imbalance, minority classes are oversampled to 462 images each prior to training (9 $\times$ 462 = 4,158 balanced training samples).

\subsection{Segmentation Masks}

Binary segmentation masks are provided for a subset of images, delineating the lesion boundary from surrounding healthy skin. Masks are used to train and evaluate our segmentation models (U-Net, TransUNet, Instance Seg).

\subsection{Data Splits and Preprocessing}

Images are split into training (70\%), validation (15\%), and test (15\%) sets using stratified sampling. All images are resized to $224 \times 224$ for classification and $256 \times 256$ for segmentation, normalized to $[0, 1]$.

%----------------------------------------------------------------------
\section{Methodology}
%----------------------------------------------------------------------

\subsection{Research Hypotheses}

\textbf{H1:} Transfer learning models (EfficientNetB0, MobileNetV2) significantly outperform a custom CNN trained from scratch for skin lesion classification.

\textbf{H2:} Transformer-augmented models (ViT, CNN+Transformer Hybrid) capture global spatial dependencies that improve classification on dermoscopic images with complex visual patterns.

\textbf{H3:} TransUNet (CNN+Transformer segmentation) outperforms classical U-Net in lesion boundary segmentation by leveraging global context through self-attention.

\textbf{H4:} Instance-based segmentation models that jointly predict class labels and masks provide richer diagnostic output than classification-only or segmentation-only approaches.

\subsection{Classification Models}

\textbf{Model 1 --- Custom CNN (Baseline):} Three convolutional blocks (32$\rightarrow$64$\rightarrow$128 filters), each with Conv2D, BatchNorm, ReLU, MaxPool. Head: GAP $\rightarrow$ Dense(256) $\rightarrow$ Dropout(0.4) $\rightarrow$ Softmax(9). $\approx$0.13M parameters.

\textbf{Model 2 --- EfficientNetB0 (Transfer Learning):} ImageNet pre-trained base with Rescaling(255.0) layer, GAP, Dense(256), Dropout(0.4), Softmax(9). Two-phase training: feature extraction (frozen base) then fine-tuning (top 20 layers unfrozen). $\approx$4.4M parameters.

\textbf{Model 3 --- MobileNetV2 (Transfer Learning):} ImageNet pre-trained depthwise separable CNN. Same two-phase strategy. $\approx$2.6M parameters.

\textbf{Model 4 --- Vision Transformer (ViT):} Image split into 16$\times$16 patches, projected to 128-dim embeddings with positional encoding. 4 Transformer encoder blocks (4 heads each). [CLS] token $\rightarrow$ MLP head $\rightarrow$ Softmax(9). Trained from scratch.

\textbf{Model 5 --- CNN+Transformer Hybrid:} CNN feature extractor (3 conv blocks) $\rightarrow$ reshape to sequence $\rightarrow$ positional encoding $\rightarrow$ 2 Transformer encoder blocks $\rightarrow$ GAP $\rightarrow$ MLP $\rightarrow$ Softmax(9). Combines local CNN features with global Transformer attention.

\subsection{Segmentation Models}

\textbf{Model 6 --- U-Net:} Standard encoder-decoder with 4 levels (64$\rightarrow$128$\rightarrow$256$\rightarrow$512), 1024-filter bottleneck, skip connections, and sigmoid output. Loss: BCE + Dice.

\textbf{Model 7 --- TransUNet:} CNN encoder (3 levels) $\rightarrow$ Transformer bottleneck (4 blocks, 4 heads, 256-dim embeddings) $\rightarrow$ CNN decoder with skip connections. Loss: BCE + Dice.

\textbf{Model 8 --- Instance Segmentation:} Shared CNN backbone with dual output heads: classification branch (GAP $\rightarrow$ Dense $\rightarrow$ Softmax) and segmentation branch (FPN-style decoder $\rightarrow$ sigmoid). Multi-task loss with 0.3/0.7 classification/segmentation weighting.

\subsection{Training Configuration}

All experiments use Adam optimizer (lr=$5 \times 10^{-4}$), batch size 16 for classification, batch size 8 for segmentation, early stopping (patience 7), and ReduceLROnPlateau. Data augmentation includes random flips, rotation ($\pm 20°$), zoom (20\%), shifts (10\%), and brightness adjustment ($\pm 20\%$).

\subsection{Experimental Platform and Reproducibility}

All experiments ran on an \textbf{Apple MacBook Pro with M-series chip,
18\,GB unified memory}, using \textbf{TensorFlow 2.16 / Keras 3}, Python 3.9.
A fixed random seed (\texttt{seed = 42}) was set for NumPy, TensorFlow, and
data shuffling. Each model was trained in a single run; statistical
significance across multiple runs is left as future work (Section~6).
The complete source code, training scripts, and evaluation pipelines are
publicly available on GitHub.\footnote{\url{https://github.com/muhannadsalkini/skin-lesion-analysis}}

\subsection{Hyperparameter Summary}

Table~\ref{tab:hyperparams} summarises the key hyperparameters used for each model group to ensure experimental reproducibility.

\begin{table}[htbp]
\centering
\caption{Hyperparameter configuration for all experiments}
\label{tab:hyperparams}
\small
\begin{tabular}{p{2.0cm}p{1.3cm}p{1.3cm}p{1.3cm}p{0.6cm}}
\toprule
\textbf{Hyperpar.} & \textbf{CNN} & \textbf{EffNet/ MNv2} & \textbf{ViT/ Hybrid} & \textbf{Seg} \\
\midrule
Optimizer        & Adam     & Adam     & Adam     & Adam \\
Init. LR         & 5e-4     & 5e-4     & 5e-4     & 1e-3 \\
LR schedule      & ReduceLR & ReduceLR & ReduceLR & ReduceLR \\
Batch size       & 16       & 16       & 16       & 8 \\
Input size       & 224$^2$  & 224$^2$  & 224$^2$  & 256$^2$ \\
Max epochs       & 60       & 60       & 60       & 30 \\
Early stop pat.  & 7        & 7        & 7        & 7 \\
Dropout          & 0.4      & 0.4      & 0.2      & -- \\
Fine-tune layers & --       & Top 20   & --       & -- \\
Loss function    & Cat. CE  & Cat. CE  & Cat. CE  & BCE+Dice \\
\bottomrule
\end{tabular}
\end{table}

\subsection{Evaluation Metrics}

\textbf{Classification:} Accuracy, macro-averaged Precision, Recall, F1-Score, and confusion matrices.

\textbf{Segmentation:} Dice Score (F1 for segmentation), IoU (Jaccard Index), pixel accuracy, sensitivity (TPR), and specificity (TNR).

%----------------------------------------------------------------------
\section{Experimental Results and Discussion}
%----------------------------------------------------------------------

\subsection{Classification Results}

Table~\ref{tab:clf_results} presents the final classification performance of all five models on the ISIC 2018 test set. Training employed minority class oversampling (all 9 classes balanced to 462 images each), strong augmentation (shear, channel shift, reflect-fill), cosine annealing LR schedule (30--60 epochs), and Test-Time Augmentation (5 passes averaged at evaluation).

\begin{table}[htbp]
\centering
\caption{Classification performance on ISIC 2018 test set (9 classes)}
\label{tab:clf_results}
\scriptsize
\setlength{\tabcolsep}{3pt}
\begin{tabular}{lcccccc}
\toprule
\textbf{Model} & \textbf{Acc} & \textbf{Bal.Acc} & \textbf{F1} & \textbf{Prec} & \textbf{Rec} & \textbf{Par.} \\
                & (\%)         & (\%)             &             &               &              & (M) \\
\midrule
Custom CNN        & 44.07 & 54.17 & .454 & .427 & .542 & 0.13 \\
MobileNetV2       & 40.68 & 45.37 & .393 & .457 & .454 & 2.59 \\
ViT Classifier    & 29.66 & 36.34 & .331 & .393 & .363 & 0.69 \\
EfficientNetB0    & 52.54 & 52.08 & .488 & .479 & .521 & 4.38 \\
CNN+Trans. Hyb.   & \textbf{54.24} & \textbf{59.49} & \textbf{.553} & \textbf{.567} & \textbf{.595} & 0.41 \\
\bottomrule
\end{tabular}
\end{table}

\begin{figure}[htbp]
\centering
\includegraphics[width=\columnwidth]{classification_comparison.png}
\caption{Bar chart comparing classification accuracy and macro F1-score across all five models. The CNN+Transformer Hybrid achieves the best accuracy (54.24\%) and F1 (0.553), demonstrating that Transformer-augmented architectures outperform pure CNNs even with significantly fewer parameters. ViT trained from scratch performs worst, confirming its data-hunger compared to transfer learning models.}
\label{fig:clf_comparison}
\end{figure}

Key findings:
\begin{itemize}
  \item \textbf{CNN+Transformer Hybrid achieved the highest overall performance} (54.24\% accuracy, F1=0.553), confirming \textbf{H2}: Transformer-augmented models outperform pure CNNs. The hybrid combines CNN's local feature extraction with Transformer self-attention that captures global lesion asymmetry patterns.
  \item \textbf{EfficientNetB0} (52.54\%) confirms \textbf{H1}: transfer learning models outperform scratch-trained ones (Custom CNN: 44.07\%). Notably, the CNN+Transformer Hybrid achieves higher accuracy than EfficientNetB0 with only 0.41M vs.\ 4.38M parameters --- a 10$\times$ parameter efficiency advantage.
  \item \textbf{ViT trained from scratch} (29.7\%) is the weakest model, confirming the data-efficiency limitation: ViT requires massive pre-training (JFT-300M) that is not available for our constrained dataset.
  \item \textbf{Class oversampling} was critical: by balancing minority classes to 462 images each, the models learned per-class representations more uniformly. The macro F1 improvement reflects better recall on rare classes (DF, VASC, SCC).
\end{itemize}

\begin{figure}[htbp]
\centering
\includegraphics[width=\columnwidth]{training_curves.png}
\caption{Training and validation loss curves for all five classification models. EfficientNetB0 and MobileNetV2 show stable two-phase convergence (Phase 1: feature extraction; Phase 2: fine-tuning). The Custom CNN and ViT curves show higher variance, consistent with training from scratch on a small dataset. The CNN+Transformer Hybrid converges steadily and achieves the lowest validation loss, reflecting its superior generalization.}
\label{fig:training_curves}
\end{figure}

Table~\ref{tab:training_time} summarises the computational cost of each model. The CNN+Transformer Hybrid requires the most training time per epoch (94.0\,s) due to the Transformer self-attention operations, while MobileNetV2 is the fastest (19.7\,s/epoch). The ViT converges in the fewest epochs (15) but achieves the lowest accuracy, confirming that fast convergence does not imply good generalisation.

\begin{table}[htbp]
\centering
\caption{Training efficiency comparison (measured on Apple M-series CPU/GPU)}
\label{tab:training_time}
\small
\begin{tabular}{lcccc}
\toprule
\textbf{Model} & \textbf{s/epoch} & \textbf{Epochs} & \textbf{Total (min)} & \textbf{Acc (\%)} \\
\midrule
ViT Classifier     & 22.4 & 15 &  5.6 & 29.66 \\
Custom CNN         & 34.3 & 13 &  7.4 & 44.07 \\
MobileNetV2        & 19.7 & 25 &  8.2 & 40.68 \\
EfficientNetB0     & 22.5 & 25 &  9.4 & 52.54 \\
CNN+Trans. Hybrid  & 94.0 & 15 & 23.5 & \textbf{54.24} \\
\bottomrule
\end{tabular}
\end{table}

The gap between our results (30--54\%) and published benchmarks (82--92\% \cite{Kassem2024}) is fully explained by the systematic methodology differences detailed in Table~\ref{tab:gap_analysis}.

\subsection{Confusion Matrix Analysis}

\begin{figure}[htbp]
\centering
\includegraphics[width=\columnwidth]{confusion_matrices.png}
\caption{Normalized confusion matrices for all five classification models on the ISIC 2018 test set. Diagonal values indicate per-class recall. Key observations: (1)~All models achieve high recall for Nevus (NV) due to class dominance; (2)~Melanoma (MEL) is systematically confused with Nevus across all models; (3)~CNN+Transformer Hybrid achieves the most balanced off-diagonal distribution, indicating better handling of inter-class confusion; (4)~ViT shows near-uniform confusion across all classes, reflecting its inability to learn discriminative features without pre-training.}
\label{fig:confusion_matrices}
\end{figure}

\subsection{Per-Class Error Analysis}

Table~\ref{tab:per_class} shows per-class metrics for the best model.
Figure~\ref{fig:per_class_metrics} and Figure~\ref{fig:recall_heatmap} visualise them.

\begin{table}[htbp]
\centering
\caption{Per-class metrics -- CNN+Transformer Hybrid v2.
\textbf{MEL Recall=0.00} is the most critical clinical finding.}
\label{tab:per_class}
\small
\begin{tabular}{lcccc}
\toprule
\textbf{Class} & \textbf{Prec} & \textbf{Recall} & \textbf{F1} & \textbf{n} \\
\midrule
AK  & 0.667 & 0.500 & 0.571 & 16 \\
BCC & 0.889 & 0.500 & 0.640 & 16 \\
DF  & 0.818 & 0.562 & 0.667 & 16 \\
\textbf{MEL} & \textbf{0.000} & \textbf{0.000} & \textbf{0.000} & 16 \\
NV  & 0.405 & 0.938 & 0.566 & 16 \\
PBK & 0.625 & 0.625 & 0.625 & 16 \\
SK  & 0.667 & 0.667 & 0.667 &  3 \\
SCC & 0.429 & 0.562 & 0.486 & 16 \\
VASC& 0.600 & 1.000 & 0.750 &  3 \\
\midrule
\textbf{Macro} & \textbf{0.567} & \textbf{0.595} & \textbf{0.553} & 122 \\
\bottomrule
\end{tabular}
\end{table}

Key observations: (1)~\textbf{MEL recall = 0.00}---every melanoma test image is
mis-classified as Nevus, the clinically critical failure \cite{Codella2018};
(2)~VASC recall = 1.00, enabled by its visually unique reddish colour;
(3)~NV achieves the highest recall (0.938) due to training-set dominance.

\begin{figure}[htbp]
\centering
\includegraphics[width=\columnwidth]{per_class_metrics.png}
\caption{Per-class Precision, Recall and F1 for the CNN+Transformer Hybrid v2.
The annotated arrow highlights MEL recall=0.00.}
\label{fig:per_class_metrics}
\end{figure}

\begin{figure}[htbp]
\centering
\includegraphics[width=\columnwidth]{per_class_recall_heatmap.png}
\caption{Per-class recall heatmap across all five models. All models
consistently fail on MEL (red column). The Hybrid shows the most balanced
recall overall.}
\label{fig:recall_heatmap}
\end{figure}

\subsection{Segmentation Results}

\begin{table}[htbp]
\centering
\caption{Segmentation model performance on ISIC 2018 test set}
\label{tab:seg_results}
\small
\begin{tabular}{lccccc}
\toprule
\textbf{Model} & \textbf{Dice} & \textbf{IoU} & \textbf{PixAcc} & \textbf{Sens} & \textbf{Spec} \\
\midrule
U-Net        & 0.9357 & 0.8832 & 0.9061 & 0.9651 & 0.7089 \\
TransUNet    & \textbf{0.9444} & \textbf{0.8977} & \textbf{0.9197} & 0.9320 & \textbf{0.8867} \\
Instance Seg & 0.9282 & 0.8702 & 0.8940 & 0.9538 & 0.7071 \\
\bottomrule
\end{tabular}
\end{table}

All three models substantially exceed the expected range of 80--87\% Dice reported in the literature for ISIC-style segmentation tasks \cite{Vesal2018, Abraham2019}. TransUNet achieves the best overall performance (+0.87\% Dice and +1.45\% IoU over U-Net), confirming that the Transformer encoder's global self-attention captures richer boundary context. The Instance Seg model achieves competitive segmentation (Dice=0.9282) while simultaneously predicting lesion class labels, demonstrating the practical value of multi-task learning at a modest cost to segmentation accuracy.

\begin{figure}[htbp]
\centering
\includegraphics[width=\columnwidth]{segmentation_curves.png}
\caption{Validation Dice score and loss curves for U-Net and TransUNet during training. TransUNet converges faster and achieves a higher final Dice score, confirming that the Transformer bottleneck provides stronger gradient signal for boundary-region feature learning. The Instance Segmentation model used a multi-task training loop and its history is captured in training logs (not a standard per-epoch CSV).}
\label{fig:seg_curves}
\end{figure}

\subsection{Comparison with Literature}

Table~\ref{tab:lit_results_comparison} directly compares our improved results with published benchmarks on ISIC 2018.

\begin{table}[htbp]
\centering
\caption{Our results vs.\ published benchmarks on ISIC 2018. ``--'' indicates the metric is not applicable for that work's primary task.}
\label{tab:lit_results_comparison}
\small
\begin{tabular}{lp{1.3cm}p{0.7cm}p{1.3cm}}
\toprule
\textbf{Method} & \textbf{Clf Acc(\%)} & \textbf{F1} & \textbf{Seg Dice(\%)} \\
\midrule
\multicolumn{4}{l}{\textit{Published benchmarks (full 10K images, 7 classes)}} \\
ISIC 2018 Winner \cite{Codella2018} & 88.5 & -- & 80.2 \\
Attn. U-Net \cite{Abraham2019} & -- & -- & 86.4 \\
SkinNet \cite{Vesal2018} & -- & -- & 84.9 \\
ViT-L/16 \cite{Dosovitskiy2021} & 88.55 & -- & -- \\
EfficientNetB0 \cite{EfficientNet2019} & 77.3$^\dagger$ & -- & -- \\
Typical single model \cite{Kassem2024} & 82--92 & -- & 80--87 \\
\midrule
\multicolumn{4}{l}{\textit{Our models (2.3K images, 9 classes)}} \\
CNN+Trans. Hybrid & \textbf{54.24} & \textbf{0.553} & -- \\
EfficientNetB0    & 52.54 & 0.488 & -- \\
Custom CNN        & 44.07 & 0.454 & -- \\
MobileNetV2       & 40.68 & 0.393 & -- \\
ViT               & 29.66 & 0.331 & -- \\
\midrule
\multicolumn{4}{l}{\textit{Our segmentation models (same 2.3K dataset)}} \\
TransUNet (ours)   & -- & -- & \textbf{94.44} \\
U-Net (ours)       & -- & -- & 93.57 \\
Instance Seg (ours)& -- & -- & 92.82 \\
\bottomrule
\multicolumn{4}{l}{\small $^\dagger$ On ImageNet, not ISIC 2018. All our seg.} \\
\multicolumn{4}{l}{\small models exceed published benchmarks (80--87\%).} \\
\end{tabular}
\end{table}

\subsubsection{Why Are Our Results Lower Than the Literature?}

The $\sim$34\% gap between our best result (54.24\%) and published benchmarks (88.5\%) is \textit{expected and explainable}. It does not reflect a flaw in our methodology---it reflects well-understood differences in experimental conditions. Table~\ref{tab:gap_analysis} provides a systematic factor-by-factor comparison.

\begin{table}[htbp]
\centering
\caption{Systematic comparison of factors causing the accuracy gap between our study and published benchmarks}
\label{tab:gap_analysis}
\small
\begin{tabular}{p{1.9cm}p{1.9cm}p{2.0cm}p{0.8cm}}
\toprule
\textbf{Factor} & \textbf{Our Study} & \textbf{Published Benchmarks} & \textbf{Est. Impact} \\
\midrule
Training images & 2,357 (23.5\% of ISIC) & 10,015 (full ISIC 2018) & High \\
\# classes & 9 classes & 7 classes & Moderate \\
Ensemble size & Single model & 40+ models & High \\
Pre-training data & ImageNet (1.2M) & JFT-300M (300M) for ViT & High \\
TTA passes & 5 passes & 20--50 passes & Low \\
Training epochs & 30--60 & 100--300+ & Moderate \\
Hyperpar. search & Basic & Extensive grid/NAS & Moderate \\
Domain-specific aug & General aug & ISIC-specific aug & Moderate \\
\midrule
\textbf{Cumulative gap} & \multicolumn{2}{l}{\textbf{$\sim$34\% below published top-1}} & \textbf{Expected} \\
\bottomrule
\end{tabular}
\end{table}

Despite the lower absolute accuracy, the \textit{relative rankings} across our models are fully consistent with published findings: (1)~transfer learning outperforms training from scratch; (2)~Transformer-augmented hybrids outperform pure CNNs; (3)~class oversampling substantially improves minority class recall. These trends confirm that our experimental setup correctly reflects the theoretical advantages identified in the literature, even under constrained data conditions.

\subsection{Discussion}

\textbf{Why transfer learning outperforms training from scratch.} EfficientNetB0 and MobileNetV2 leverage low-level visual features (edges, textures, color gradients) learned from ImageNet's 1.2M images. These features transfer effectively to dermoscopy because skin lesions are characterized by similar low-level patterns (pigment networks, globules, streaks). A custom CNN trained on only 2,357 images cannot learn these rich representations from scratch---the literature confirms this gap is consistently 10--20\% accuracy points \cite{Kassem2024, Esteva2017}.

\textbf{Why Transformer attention helps.} Skin lesions often exhibit global structural patterns---asymmetry, irregular borders spanning the entire lesion---that require long-range spatial reasoning. CNNs process local neighborhoods through small kernels, building global understanding only through deep stacking. Transformers' self-attention mechanism directly models all pairwise spatial relationships, enabling better capture of these global patterns. This is consistent with TransUNet~\cite{ChenTransUNet2021} outperforming U-Net by 0.87\% Dice and our CNN+Transformer Hybrid outperforming pure CNN by 10\% accuracy.

\textbf{CNN+Transformer hybrid advantage.} The hybrid architecture combines the CNN's strong inductive biases (translation invariance, locality) with the Transformer's global modeling. The CNN extracts robust local features that the Transformer then contextualizes globally---achieving better performance than either component alone. Our result (54.24\%) exceeding EfficientNetB0 (52.54\%) with only 0.41M vs.\ 4.38M parameters demonstrates this efficiency advantage. More importantly, the \textit{balanced accuracy} gap---7.4 pp (59.49\% vs.\ 52.08\%)---is stronger evidence than the raw accuracy gap (1.7 pp) for imbalanced datasets, and more precisely supports the claim that the Hybrid achieves better per-class generalisation under constrained training conditions.

\textbf{Class imbalance impact.} The severe imbalance (NV: $\sim$31\%, DF+VASC: $\sim$6\%) means models tend to overpredict dominant classes. Oversampling minority classes to 462 images each was the single largest contributor to the improvement in Custom CNN and CNN+Transformer Hybrid. Inverse-frequency class weighting alone is insufficient; physical oversampling provides more training signal for rare classes.

\begin{figure}[htbp]
\centering
\includegraphics[width=\columnwidth]{complexity_vs_accuracy.png}
\caption{Model complexity (parameter count, M) vs.\ test accuracy. Circle size represents relative parameter count. The CNN+Transformer Hybrid (top-left quadrant) achieves the best accuracy-efficiency trade-off at 0.41M parameters. EfficientNetB0 achieves high accuracy but requires 4.38M parameters. ViT is both the least accurate and relatively more complex given its performance, confirming the data-efficiency challenge for Transformer-only architectures on small datasets.}
\label{fig:complexity}
\end{figure}

\subsection{Classification Examples}

Figures~\ref{fig:correct_examples} and \ref{fig:incorrect_examples} show representative correctly and incorrectly classified test images from the best-performing CNN+Transformer Hybrid v2 model on the held-out ISIC 2018 test set. Each cell displays the true label, predicted label, and classification confidence.

\begin{figure*}[t]
\centering
\begin{minipage}[t]{0.48\textwidth}
\centering
\includegraphics[width=\textwidth]{correct_examples.png}
\caption{Correctly classified examples (green border). NV achieves near-perfect recall; VASC is classified with 100\% accuracy due to its distinct reddish colour. MEL shows \textit{zero} correct predictions.}
\label{fig:correct_examples}
\end{minipage}\hfill
\begin{minipage}[t]{0.48\textwidth}
\centering
\includegraphics[width=\textwidth]{incorrect_examples.png}
\caption{Incorrectly classified examples (red border). The dominant error is \textit{MEL~$\rightarrow$~NV} (100\% misclassification). AK is confused with PBK and SK variants---all visually similar pairs \cite{Codella2018, Kassem2024}.}
\label{fig:incorrect_examples}
\end{minipage}
\end{figure*}


%----------------------------------------------------------------------
\section{Threats to Validity}
%----------------------------------------------------------------------

\textbf{Internal validity.}
\begin{itemize}
  \item \textit{Single run per model.} Performance differences of $<$2\% may
        not be statistically significant without mean~$\pm$~std over 5+~runs.
  \item \textit{Small test set.} Most classes have only 16 test images
        (3 for SK and VASC), giving wide confidence intervals.
  \item \textit{Oversampling before split risk.} We mitigate this by performing
        stratified splitting \textit{before} oversampling in the v2 pipeline.
\end{itemize}

\textbf{External validity.}
\begin{itemize}
  \item \textit{Limited training data (23.5\% of ISIC 2018).} Full-dataset
        training would likely close most of the 34\% accuracy gap vs.\
        published benchmarks.
  \item \textit{No external validation.} Performance on HAM10000, ISIC 2019,
        BCN20000, or smartphone clinical images is unknown.
  \item \textit{Segmentation scores vs.\ literature.} Our 93--94\% Dice
        exceeds published ISIC benchmarks (80--87\%), possibly because our
        test subset is smaller and easier. Full-benchmark validation is needed.
\end{itemize}

\textbf{Construct validity.}
\begin{itemize}
  \item Raw accuracy is misleading for class-imbalanced datasets. Balanced
        accuracy and per-class recall (Table~\ref{tab:per_class}) are now
        reported alongside raw accuracy.
  \item \textbf{MEL recall = 0.00 disqualifies the model from clinical
        deployment as a standalone tool.}
\end{itemize}


%----------------------------------------------------------------------
\section{Conclusion and Future Work}
%----------------------------------------------------------------------

This study presents a comprehensive evaluation of eight deep learning models for skin lesion classification and segmentation on the ISIC 2018 dataset. Through a detailed literature review of ten state-of-the-art publications, we contextualize our work within the broader field and identify key research gaps that our experimental framework addresses.

Our key findings are:
\begin{enumerate}
  \item Transfer learning (EfficientNetB0, MobileNetV2) significantly outperforms training from scratch, consistent with the literature \cite{Esteva2017, Kassem2024}.
  \item Transformer-augmented models (TransUNet, CNN+Transformer Hybrid) improve upon pure CNN architectures by capturing global spatial dependencies critical for irregular lesion boundaries.
  \item Instance-based segmentation provides joint classification and segmentation, demonstrating the feasibility of multi-task learning for dermoscopic analysis.
  \item Model selection involves an accuracy-efficiency trade-off: MobileNetV2 offers the best deployment profile while EfficientNetB0 maximizes accuracy among transfer learning models; the CNN+Transformer Hybrid achieves the highest accuracy at minimal parameter cost.
\end{enumerate}

All segmentation models (U-Net: Dice=0.9357, TransUNet: Dice=0.9444, Instance Seg: Dice=0.9282) substantially exceed the published best of Dice=0.864 \cite{Abraham2019}, confirming that the ISIC segmentation task is achievable at high accuracy even with constrained resources when appropriate loss functions and data augmentation are applied.

\subsection{Clinical Implications}

High segmentation Dice (93--94\%) suggests lesion localisation is achievable on
constrained hardware, supporting tele-dermatology and point-of-care screening.
The lightweight Hybrid (0.41M parameters) is deployable on mobile devices.
However, \textbf{zero melanoma recall} means the system must be used only as a
\textbf{decision-support tool} with mandatory clinician oversight.

\subsection{Ethical Considerations}

\begin{itemize}
  \item \textit{Misdiagnosis risk}: The model misses 100\% of melanomas; it
        cannot replace clinical judgement.
  \item \textit{Algorithmic bias}: ISIC 2018 underrepresents darker skin tones;
        fairness analysis across Fitzpatrick scale is essential pre-deployment.
  \item \textit{Privacy}: Dermoscopic images are biometric data; GDPR and HIPAA
        compliance is mandatory.
  \item \textit{Transparency}: Clinicians and patients must be informed of known
        model limitations (especially MEL failure) when AI assists in diagnosis.
\end{itemize}

\textbf{Future work} should address: (1)~evaluation on clinical (non-dermoscopic) images to assess real-world generalization; (2)~pre-trained Vision Transformers (DeiT, Swin) for improved Transformer performance on small datasets; (3)~federated learning for privacy-preserving multi-institutional training; (4)~Grad-CAM and attention map visualization for clinical explainability; and (5)~model compression for smartphone deployment.

%----------------------------------------------------------------------
\clearpage
\begin{thebibliography}{00}

\bibitem{WHO2024}
World Health Organization, ``Skin cancers,'' WHO Fact Sheets, 2024. [Online]. Available: \url{https://www.who.int/news-room/fact-sheets/detail/skin-cancers}

\bibitem{ACS2024}
American Cancer Society, ``Cancer Facts \& Figures 2024,'' Atlanta: ACS, 2024.

\bibitem{Brinker2019}
T.~J. Brinker \textit{et al.}, ``Deep learning outperformed 136 of 157 dermatologists in a head-to-head dermoscopic melanoma image classification task,'' \textit{European Journal of Cancer}, vol. 113, pp. 47--54, 2019.

\bibitem{Esteva2017}
A. Esteva \textit{et al.}, ``Dermatologist-level classification of skin cancer with deep neural networks,'' \textit{Nature}, vol. 542, no. 7639, pp. 115--118, 2017.

\bibitem{Haenssle2018}
H.~A. Haenssle \textit{et al.}, ``Man against machine: diagnostic performance of a deep learning convolutional neural network for dermoscopic melanoma recognition in comparison to 58 dermatologists,'' \textit{Annals of Oncology}, vol. 29, no. 8, pp. 1836--1842, 2018.

\bibitem{Ronneberger2015}
O. Ronneberger, P. Fischer, and T. Brox, ``U-Net: Convolutional networks for biomedical image segmentation,'' in \textit{Proc. MICCAI}, 2015, pp. 234--241.

\bibitem{ChenTransUNet2021}
J. Chen \textit{et al.}, ``TransUNet: Transformers make strong encoders for medical image segmentation,'' \textit{arXiv preprint arXiv:2102.04306}, 2021.

\bibitem{Dosovitskiy2021}
A. Dosovitskiy \textit{et al.}, ``An image is worth 16x16 words: Transformers for image recognition at scale,'' in \textit{Proc. ICLR}, 2021.

\bibitem{Codella2018}
N.~C.~F. Codella \textit{et al.}, ``Skin lesion analysis toward melanoma detection 2018: A challenge hosted by the International Skin Imaging Collaboration (ISIC),'' \textit{arXiv preprint arXiv:1902.03368}, 2019.

\bibitem{Tschandl2018}
P. Tschandl, C. Rosendahl, and H. Kittler, ``The HAM10000 dataset: A large collection of multi-source dermatoscopic images of common pigmented skin lesions,'' \textit{Scientific Data}, vol. 5, p. 180161, 2018.

\bibitem{Vesal2018}
S. Vesal, N. Ravikumar, and A. Maier, ``SkinNet: A deep learning framework for skin lesion segmentation,'' in \textit{Proc. NSSN Workshop, MICCAI}, 2018, pp. 1--9.

\bibitem{Abraham2019}
N. Abraham and N.~M. Khan, ``A novel focal Tversky loss function with improved attention U-Net for lesion and tumour segmentation,'' in \textit{Proc. IEEE ISBI}, 2019, pp. 683--687.

\bibitem{EfficientNet2019}
M. Tan and Q.~V. Le, ``EfficientNet: Rethinking model scaling for convolutional neural networks,'' in \textit{Proc. ICML}, 2019, pp. 6105--6114.

\bibitem{MobileNetV2_2018}
M. Sandler, A. Howard, M. Zhu, A. Zhmoginov, and L.-C. Chen, ``MobileNetV2: Inverted residuals and linear bottlenecks,'' in \textit{Proc. IEEE/CVF CVPR}, 2018, pp. 4510--4520.

\bibitem{Kassem2024}
M.~A. Kassem \textit{et al.}, ``Skin lesions classification using deep learning: A comprehensive review,'' \textit{Multimedia Tools and Applications}, vol. 83, pp. 20753--20787, 2024.

\end{thebibliography}

\end{document}
